% Exchange the input and encoded unit positions; preserves the 110 milestone.
\section{Putting the input in the unit position}\label{sec:input-unit}

The supplied input can occupy the constant coefficient of the code.
Exchanging its position with that of the encoded unit changes
$\mathcal C=xB+g$ to $\mathcal C=x+g$, removing one multiplication.
The main proof obligation is to recover the Pell conclusions without
assuming $\mathcal C>B$. The geometric equation provides the weaker
bound $B\le q^2$, which suffices.

\subsection{Remapping the homogeneous code}

Retain the homogeneous rows, guard $u\delta-x^2$, special target
$\delta^2$, and all numerical layout constants of
Section~\ref{sec:unit-digit}. Give $x$ weight zero and $\delta$
weight $v_0=1$. The original coordinates retain weights
$v_1,\ldots,v_{58}$, and $u$ retains weight $v_{59}$.
The mask becomes
\[
 \ell_0(B)=\sum_{i=0}^{59}B^{v_i}
                    +\sum_{j=0}^{1831}B^{t_j}.
\]
It excludes position zero, which now belongs to the input.
There are still 1892 non-input encoded coordinates and 1893 coordinates
including the input. The same bound $1900b$ applies to their sum.

Recompute each complementary coefficient position using the new
monomial weights. The ordinary coefficient rule is still $2a/c$.
The special row has coefficient $1$ at $t_{1831}-2$, since
$\delta^2$ now has weight two. The guard's two monomial weights are
$v_{59}+1$ and zero. All 1891 homogeneous monomial weights remain
distinct: exchanging which coordinate has weight zero does not affect
the number of positive-weight factors or the Sidon argument.
All support-separation bounds consequently remain valid.

With the remapped signed polynomial $\mathcal D$, define
\[
 e_0(B)=z\sum_{h=0}^{K-1}B^h+\mathcal D(B),
 \qquad V=\ell_0(Z)+e_0(Z)Z^L,
\]
and use the same admissibility bound \eqref{eq:unit-H}.
These indices are again fixed independently of the input value.
Make the sole arithmetic replacement
\[
 \mathcal C=x+g
\]
in the 110-operation system. In particular $Y=\ell+eq$,
$Y+\mathcal C^2+\alpha=q^2$, $\sigma=S_3>0$, the mask
$(B-4)\ell$, and all Pell equations remain unchanged.

\subsection{The weaker bounds before the Pell step}

The positive bound gives
\[
 0<\ell<Y<q^2,\qquad e<q,\qquad
 \mathcal C<q,\qquad g<q.
\]
The geometric equation and $\lambda\ge1$ imply
\begin{equation}\label{eq:input-weaker-radix}
 q^2-1=\lambda(B-1)\ge B-1,
 \qquad B\le q^2\le n=q^8.
\end{equation}
Since $B=Hb^2$ and $H>16$, we already have $q>4b\ge8$.
Put $\mathcal N=(B-Z)\lambda$. Then
\[
 b<B-Z\le\mathcal N<q^2-1,
 \qquad B\lambda<2q^2.
\]
No inequality $\lambda>q$ is needed. The imposed positive register
and its elementary upper bound give
\[
 0<S_3<2q^3+2q^2(1+q)<5q^3<q^4.
\]
Thus $S=g+q^2(Y+q^2S_3)$ satisfies $0<S<n$.

For completeness, the possibly negative first raw mask still has
an exact normalized packing. Set
\[
 d_0=\left\lfloor\frac{(b-1)\ell}{q^2}\right\rfloor,
 \qquad \varepsilon_0=(b-1)\ell-d_0q^2.
\]
Then $0\le d_0<b$, $0\le\varepsilon_0<q^2$, and
\[
 T=(q^2-1-\varepsilon_0)
       +q^2(\mathcal N-d_0)+q^4(B-4)\ell.
\]
The first two blocks belong to $[0,q^2)$, since
$\mathcal N>b>d_0$. By \eqref{eq:input-weaker-radix},
the last block satisfies
$0<(B-4)\ell<q^4$. Hence $0\le T<n$, and
\[
 r=S(n^2-n)+(T+1)(n^2-1)\ge n^2-1\ge n.
\]
This establishes all packing bounds before decoding or any conclusion
that $B<q$.

\subsection{The Pell bootstrap with \texorpdfstring{$B\le n$}{B at most n}}

Here is the dependency check for the relaxed radix comparison.
The information now available is
\begin{equation}\label{eq:input-pell-initial}
 3<3L\le B\le n\le r,\qquad q,b\le n,
 \qquad n,r\ge8.
\end{equation}
Moreover $B$ is even because the fixed parameter $H$ is an even
power of two, independently of any assertion about $b$.

Write $N=n$, $R=r$, $U=wN^2$, $Y_P=sN^2$,
$M=RY_P$, $a=M(U+1)$, $X_P=2UM^2$, $P=X_P+1$, and
$J=2R+1$. We have $U,Y_P\ge64$ and $M\ge512$.
The index congruence gives $k\ge R+2$, and the positive interval
$Y_P<c/k<Y_P+1$ gives $c>J$. The signed Pell block therefore
gives $c=\psi_a(J)$ and $d=\chi_a(J)$.
The growth and index-congruence argument of
Section~\ref{sec:shifted-pell} gives $k=\psi_P(R+1)$ and
\[
 \xi\left(1-\frac{R}{a}-\frac{R}{X_P}\right)
       <\frac ck<\xi,
 \qquad \xi=\frac{(U+1)^{2R}}{U^R},
 \qquad \frac{R}{a}+\frac{R}{X_P}
       <\frac{2}{Y_P(U+1)}<\frac12.
\]
These arguments use the numerical lower bounds just stated, not
a comparison between $B$ and $q$. The interval consequently gives
$Y_P>U^{R-1}$ and $a>RU^R$.

The base-two exponent argument has the same strict size bounds:
\[
 2^{3J}=8\cdot8^{2R}\le8N^{2R}\le RU^R<a,
 \qquad (bw)^3\le U^R<a.
\]
Here $b\le N$ gives $bw\le U$.
Equation $E14$ and the exponent lemma yield $bw=2^J$;
thus $b$ is a power of two and $U>4\cdot2^{2R}$.

Next \eqref{eq:input-pell-initial} gives $0<L<J<a$.
The positive gap $c>\kappa$ and equations $E19,E20$ give
$\kappa=\psi_a(L)$ and $\mu=\chi_a(L)$. The second exponent
argument now uses
\begin{equation}\label{eq:input-exponent-bounds}
 B^{3L}\le B^B\le N^N\le U^R<a,
 \qquad q^3\le N^N\le U^R<a.
\end{equation}
These follow from $3L\le B\le N$, $q\le N$, and $R\ge N$.
Also $a>N\ge B$, so the modulus in $E18$ has its required
positive sign. Equation $E18$ therefore gives $q=B^L$.
Only now do we infer $B<q$.

As $B$ was already even, $q$ and $N$ are even. To make the final
rounding dependency explicit, the exact-index estimates give
\[
 0<\xi-c/k<6/(U+1)<1/4,
 \qquad 0<\xi-\lfloor\xi\rfloor<1/4,
 \qquad \lfloor\xi\rfloor\equiv\binom{2R}{R}\pmod U.
\]
The interval makes $|Y_P-\lfloor\xi\rfloor|<2$.
Both integers are even, so they are equal. Since $Y_P=sN^2$
and $N^2\mid U$, we obtain
$N^2\mid\binom{2R}{R}$. The binomial criterion and the packing
lemma now give $\tau_2(S,T)=0$. Thus the complete Pell conclusion
was recovered using \eqref{eq:input-weaker-radix}, without an
unproved early inequality $B<q$.

\subsection{Decoding and positive necessity}

The subsequent first-block borrow argument applies unchanged.
It proves $d_0=0$ and decodes
\[
 \ell=\ell_0(B)+mq,\qquad e=e_0(B)-m,
 \qquad 0\le m<e_0(B).
\]
Since $g<q$, the first mask restricts its digits to the positive
positions of $\ell_0$. Its unit digit is zero. The inequality
$x<b<B$ then ensures that $\mathcal C=x+g$ has exactly the supplied
input as its unit digit, without a carry; its digit at position one
is $\delta$.

The sum of the coefficients of this digit polynomial is positive
and less than $1900b$. All degree and coefficient estimates in
Section~\ref{sec:unit-digit} still hold. The same high-digit window
bounds the digits of $(B-4)m$ by $Z\mathcal C(1)^2-1$; evaluating
its digit polynomial at four forces $m=0$. The canonical relaxed
mask then forces every ordinary homogeneous residual and the guard
to vanish, while its special target gives $\delta^2\in\{0,1\}$.
The guard and $x>0$ give $\delta=1$, so all original residuals vanish.

Conversely, choose genuine original coordinates, set $\delta=1$
and $u=x^2$, and choose a power of two $b$ greater than the input
and every encoded coordinate. The digit $\delta$ at position one
ensures $g\ge B>0$. The canonical code has the same degree and
coefficient bounds as before, giving positive $\alpha$, $\sigma$,
and the packed quotient. All masks hold. In this necessity direction
the intended code even satisfies $\mathcal C>B$, so the preceding
Pell necessity constructions retain their stronger original bounds.

\begin{theorem}\label{thm:best109}
For each r.e.\ set, the remapped effective admissible index gives
a membership-equivalent system in 33 positive unknowns and 21 equations
with a straight-line certificate of \textbf{109 arithmetic instructions}:
60 multiplications and 49 additions. Generating the fixed numerals
from $1$ adds seven instructions, giving 116 operations.
\end{theorem}
\begin{proof}
The preceding arguments establish equivalence with all domains
preserved. Replacing $xB+g$ by $x+g$ removes exactly one multiplication
from the 110-operation schedule. No other arithmetic instruction,
unknown, or equality is added. The complete primitive list and all
21 source residual identities, including the triangular corrections,
are verified by \file{round9\_1980\_certificate.py}.
\end{proof}
